% Source-faithful assembly: Mumford, Abelian Varieties, Chapter IV,
% Horizon transcription pages 0237--0246.

\section{Theta groups and descent}
\label{mumford-IV23-continuation-source}
\dcref{M:ch4:IV.23}

\begin{proof}[Continuation of the theta-group theorem]
\label{mumford-proof-line-bundle-theta-group-continuation}
If two automorphisms of $S\times L$ covering the same translation have the
same image in $\mathcal G(L)$, their quotient is multiplication by a unit on
$S\times X$ restricting to $1$ over the chosen point. Since
$H^0(S\times X,\mathcal O^*)=H^0(S,\mathcal O_S^*)$, it is the identity.
Conversely, a point $\phi\in\mathcal G(L)(S)$ covers $f\in K(L)(S)$.
Locally on $S$, trivialize the line bundle $M$ in
$T_f^*(S\times L)\simeq(S\times L)\otimes p_1^*M$ and choose liftings
$\chi_i$ of the translations. Since $\mathcal G(L)$ is a principal
$\mathbb G_m$-bundle over $K(L)$, rescale $\chi_i$ by units so their images
are $\phi$. Injectivity makes the rescaled liftings agree on overlaps, hence
they glue. Therefore
$\underline{\operatorname{Aut}}(L/X)\simeq\mathcal G(L)$.
\source{mumford-abelian-varieties:page-0237}
\source{mumford-abelian-varieties:page-0238}
\end{proof}

\begin{definition}[Commutator pairing of a line bundle]
\label{mumford-def-line-bundle-commutator}
The theta group $\mathcal G(L)$ has a skew-symmetric bihomomorphism
\[
e^L:K(L)\times K(L)\to\mathbb G_m
\]
defined by its commutator. If $L\in\operatorname{Pic}^0X$, then $K(L)=X$,
so $e^L=1$ and $\mathcal G(L)$ is commutative.
\source{mumford-abelian-varieties:page-0238}
\end{definition}

\begin{proposition}[Functoriality of commutator pairings]
\label{mumford-prop-line-bundle-commutator-functorial}
For homomorphisms and line bundles as indicated,
\begin{enumerate}
\item $e^{f^*L}(x,y)=e^L(fx,fy)$;
\item $e^{L_1\otimes L_2}(x,y)=e^{L_1}(x,y)e^{L_2}(x,y)$;
\item algebraically equivalent line bundles have equal pairings;
\item $e^{L^n}(x,y)=e^L(x,ny)$ for $x\in K(L)$,
$y\in n_X^{-1}K(L)$;
\item if $x\in X_n$ and $y\in\phi_L^{-1}(X_n)$, then
$\bar e_n(x,\phi_L(y))=e^{L^n}(x,y)$.
\end{enumerate}
\source{mumford-abelian-varieties:page-0239}
\source{mumford-abelian-varieties:page-0240}
\source{mumford-abelian-varieties:page-0241}
\end{proposition}
\begin{proof}
Lift $f(x),f(y)$ to automorphisms in the theta group; their commutator
proves (1). Tensoring liftings proves (2), and (3) follows by writing
$L_1=L_2\otimes L_3$ with $L_3\in\operatorname{Pic}^0X$. For (4), after a
faithfully flat extension write $x=nz$; then
\[
e^L(nz,ny)=e^{n_X^*L}(z,y)=e^{L^{n^2}}(z,y)
=(e^{L^n}(z,y))^n=e^{L^n}(nz,y).
\]
For (5), write $y=nz$ and compare
$n_X^*(T_y^*L\otimes L^{-1})$ with
the line bundle $\operatorname{Hom}[n_X^*L,T_{ny}^*n_X^*L]$. The action of $x\in X_n$ on
the resulting nowhere-vanishing section is conjugation by the natural
translation isomorphism; this is simultaneously the Weil pairing and the
theta-group commutator.
\end{proof}

\begin{theorem}[Descent of line bundles]
\label{mumford-thm-theta-descent}
For an isogeny $\pi:X\to Y$ and a line bundle $L$ on $X$, line bundles $M$
on $Y$ with $\pi^*M\simeq L$ correspond naturally to homomorphisms
$\alpha:\ker\pi\to\mathcal G(L)$ lifting the inclusion into $X$.
Such an $M$ exists iff $\ker\pi\subset K(L)$ and
$e^L$ is trivial on $\ker\pi\times\ker\pi$.
\source{mumford-abelian-varieties:page-0242}
\uses{mumford-def-line-bundle-commutator}
\end{theorem}
\begin{proof}
This is the finite-group-scheme descent theorem applied to the action of
$\ker\pi$ on $L$. The lift is a splitting of the inverse image theta group.
That theta group splits iff it is commutative, equivalently iff the
commutator pairing is trivial.
\end{proof}

\begin{theorem}[Power criterion]
\label{mumford-thm-line-bundle-power-criterion}
For a line bundle $L$ and integer $n$, $L\simeq M^n$ for some $M$ iff
$K(L)\supset X_n$.
\source{mumford-abelian-varieties:page-0242}
\source{mumford-abelian-varieties:page-0243}
\end{theorem}
\begin{proof}
The forward implication follows from $\phi_{M^n}=n\phi_M$. Conversely,
$K(L)\supset X_n$ gives
$e^{L^n}(x,y)=e^L(x,ny)=1$ for $x,y\in X_n$. The descent criterion gives
$L^n\simeq n_X^*N$; then $(L\otimes N^{-n})^n\in\operatorname{Pic}^0X$.
Divisibility of $\operatorname{Pic}^0X$ supplies $P$ with
$L\otimes N^{-n}\simeq P^n$, hence $L\simeq(N\otimes P)^n$.
\end{proof}

\begin{lemma}[Orthogonal subgroup]
\label{mumford-lem-theta-orthogonal-subgroup}
For a skew bihomomorphism $e:K\times K\to\mathbb G_m$ and subgroup $H$,
$H^\perp$ is the kernel of $K\xrightarrow\gamma\widehat K\to\widehat H$;
equivalently it consists of points pairing trivially with every base change of
every point of $H$. If $L=\pi^*M$ and $H=\ker\pi$, then
$K(M)\simeq H^\perp/H$.
\source{mumford-abelian-varieties:page-0243}
\source{mumford-abelian-varieties:page-0244}
\end{lemma}
\begin{proof}
The kernel description is immediate from restriction of characters. Locally
on a base $S$, descent identifies $\pi x\in K(M)$ with an isomorphism of
$T_x^*L$ with $L$ commuting with the $H$-action. A lift $w$ of $x$ in
$\mathcal G(L)$ commutes with that action exactly when
$e^L(x,h)=1$ for all $h\in H$, which is $x\in H^\perp$.
\end{proof}

\begin{theorem}[Maximal isotropic subgroup]
\label{mumford-thm-maximal-isotropic-theta}
If $L$ is non-degenerate and $H\subset K(L)$ is maximal with
$e^L|_{H\times H}=1$, then $H=H^\perp$ and
\[
\operatorname{order}(H)^2=\operatorname{order}(K(L)).
\]
\source{mumford-abelian-varieties:page-0244}
\source{mumford-abelian-varieties:page-0245}
\end{theorem}
\begin{proof}
Put $Y=X/H$ and $L=\pi^*M$. Maximality means $M$ admits no nontrivial
further descent. If a non-degenerate bundle admits no descent of degree
$>1$, then $|\chi(M)|=1$: otherwise a prime-order subgroup of $K(M)$,
whose theta extension is commutative, would give a descent. Hence
$K(M)=0$ and $|\chi(M)|=1$. The orthogonal lemma gives $H=H^\perp$; and
Riemann--Roch gives
$|\chi(L)|=\operatorname{order}(H)$ and
$\chi(L)^2=\operatorname{order}(K(L))$.
\end{proof}

\begin{corollary}[Principal polarization and non-degeneracy]
\label{mumford-cor-principal-polarization-theta}
Every abelian variety is isogenous to a principally polarized one. If $L$ is
non-degenerate, then $\mathcal G(L)$ is a non-degenerate theta group.
\source{mumford-abelian-varieties:page-0245}
\source{mumford-abelian-varieties:page-0246}
\uses{mumford-thm-maximal-isotropic-theta}
\end{corollary}
\begin{proof}
Apply the theorem to an ample $L$ and descend to $X/H$; the descended bundle
has Euler characteristic one. For the second assertion, let $D$ be the
kernel of the commutator map $\gamma$. If $H=H^\perp$, then $D\subset H$ and
the image of $K(L)$ in $\widehat H$ has kernel $D$. Therefore
\[
\operatorname{order}K(L)\leq\operatorname{order}H\cdot
\operatorname{order}(\widehat H/D)=
\frac{\operatorname{order}(H)^2}{\operatorname{order}D}.
\]
Since $\operatorname{order}K(L)=\operatorname{order}(H)^2$, $D=1$.
\end{proof}
